The concept of fair orientations in graphs was introduced by Christodoulou, Fiat, Koutsoupias and Sgouritsa~\cite{EFXgraphs}, naturally modeling fair division scenarios in which resources are only contested by neighbors. In this model, vertices represent agents and undirected edges represent goods; edges have to be oriented towards one of their endpoints, i.e., allocated to one of their adjacent agents. Although EFX orientations (\emph{envy-free up to any good}) have been extensively studied in this setting, EF orientations (\emph{envy-free}) remain unexplored. In this work, we initiate their study primarily under the lens of parameterized complexity, presenting various tractable cases, hardness results, and parameterizations. Our results concern both simple graphs and multigraphs. Interestingly, many of our results transfer to EFX orientations, thus complementing and improving upon previous work; notably, we settle an open question regarding the structural parameterized complexity of the latter problem on graphs of polynomially-bounded valuations (Deligkas, Eiben, Goldsmith and Korchemna~\cite{ijcai/DeligkasEGK25}).
We also show that EF orientations are tractable in cases in which EFX orientations are not, particularly for binary valuations. We achieve this through a complete characterization for EF orientations with binary valuations, roughly stating that an EF orientation exists if and only if every connected component of the graph satisfies one among four graph properties. Lastly, we consider \emph{charity} in the orientation setting, establishing algorithms for finding the minimum amount of edges that have to be removed from a graph in order for EF(X) orientations to exist.


\begin{comment}
Progress:
\begin{itemize}

    \item For EF orientation, we have a DP algorithm of running time $W^{\bO(\tw)} n^{\bO(1)}$,
    where (i) the input is a simple graph and $W$ is the maximum weight,
    or (ii) the input is a multigraph and $W$ is the maximum sum of the weights of edges between the same endpoints. Sketch: keep (i) sum from forgotten edges, (ii) max sum of single edge.
    In total $W^2$ different values per vertex in the bag. The algorithm can probably be extended for EF1 and EFX orientations.
    Question: Can we get such an algorithm when the input is a multigraph and $W$ is the maximum weight of the edges?
    
    \item Towards showing that the previous algorithm for the case (i) is tight,
    we have a reduction from \RUBP~\cite{mfcs/HanakaLVY24} with polynomial weights and bounded vertex cover number.
    Maybe using some gadgets appearing in the hardness of~\cite{ijcai/DeligkasEGK25} we can show an analogous result for EFX orientations.
    
    \item towards answering the previous question, we introduced the notion of "k-balanced orientation", and showed that deciding the existence of a k-balanced orientation is FPT parameterized by k+tw. Furthermore, we showed that if a multigraph with maximum weight 2 has an orientation, then it has a 2-balanced orientation, showing that EF is FPT by tw on such graphs. 
    
    \item to complement the previous result, we show that EF on multigraphs with weight bounded by 3 is W[1]-hard parameterized by pathwidth, with a similar reduction from item 2.
    
    \item Here is a tabular that recapitulate our results for symmetric EF orientation:

    \textcolor{red}{should include something like the following table in "our contributions" section}
    
    \begin{tabular}{c|c}
        simple graphs &  multigraphs \\ \hline
        strongly NP-hard for $w=2$ & strongly NP-hard for $w=2$\\
        W[1]-hard by VC for unary weights & W[1]-hard by $|V|$ for unary weights\\
        XNLP-Hard by pw for unary weights & XNLP-Hard by pw for $w=3$\\
        XALP-Hard by tw for unary weights & XALP-Hard by tw for $w=3$\\
        FPT by tw + w & FPT by tw + W \\
         & FPT by tw if w=2 \\
        polynomial if w=1 & polynomial if w=1\\
        $\bO^* (2^k) $ algorithm where $k=\# edges\ of\ weight\ >1$ &\\
        no $\bO^* ((2-\epsilon)^k) $ algorithm under SETH & \\
         
        
    \end{tabular}
    
\end{itemize}
\end{comment}